
% 29 Mar: Clean up

\subsection{Installment 7: “Folkebladet,” April 24, 1918}

% Why the Congregation Is Not Optional [editorial title]



The controversy over Augsburg in the 1890s arose from disagreements about the concept of the church, disagreements that also shaped ministerial education. The view held by S. and O. differed greatly, especially from that of Dr. Schmidt, who in 1893 left Augsburg together with his colleagues, Professors Bøckman and Lund.

For Schmidt the outward congregation was of lesser significance. For him, what mattered above all was “the whole Christian congregation,” that is, the holy catholic Church. His concept of the church is found in Luthersk Kirkeblad, no. 41:

\begin{quote}
Christ’s true congregation is not first and foremost an outward union around certain officeholders or a certain visible organization that catches the eye, but it is the fellowship and union of true believers in the spirit, as Jesus’ true believing disciples and, each for his part, members of His spiritual body...

Thus also now, as formerly, God’s true Israel is not this or that outward flock or association of persons, a visible organization under a definite and conspicuous form of rule; rather, God’s right and genuine household are all those whom He who searches the hearts knows and loves as His own, that is, as true regenerate children of God through faith in Jesus Christ. All these true believers are members of the one body, and besides these true believers, whether they be adults or little children, there exists no member in Jesus’ body. The true and glorious unity of the Church as Jesus’ body is therefore a spiritual unity, the unity of hearts by the same spirit of faith.

Outward difference with respect to time, place, language, nation, yes, even outward ecclesiastical connection and the teaching ministry, makes no breach here in the unity of faith and spirit. All who in the wilderness of this world have learned to feel themselves bitten by the serpent, and who with the longing of faith look up to the Lamb of God, who bears the sin of the world, all these are united with Jesus by this faith that seeks salvation and are members of His body, whether moreover they be a Jew or Greek, barbarians or Scythians, whether they live in the Old or in the New Testament, whether they speak this or that tongue, yes, whether moreover they be Catholics or Protestants, or whatever kind of society they may belong to by outward confessional connection....

But if this is so, that the essential and proper unity of the Christian Church, whereby according to God’s judgment it is one single Church, is a spiritual unity, then does it perhaps follow from this that the outward and visible unity is an altogether unessential and indifferent matter?

\end{quote}

With respect to the answer that Dr. Schmidt gives, S. and O. found that the view of “the congregation” was lacking. And several years later Sv. wrote: “The core of the matter is this, which Dr. Schmidt has not yet touched: Are the New Testament’s statements about the congregation as binding upon Christians as its statements about what has been called doctrine?” This was in the years when Schmidt entered the lists on behalf of “the old church order.”

According to S. and O., therefore, the New Testament was normative also for church polity. According to Schmidt, it was a rule only for doctrine.

S. wrote: “And above all other principles in the Free Church there stands the decisive authority of the Word of God in the matter of the congregation and of church polity no less than in the question of doctrine” (Veiledning 93).

“We... have... enlightenment about ‘church polity’ in the Bible.” We need not “speculate out something ‘that suits the times’; there is light enough for the congregation in the New Testament, and it is not worth letting anyone persuade us that it does not suit our time” (105; cf. 33).

“The very first principle of the Free Church regarding church polity is this, that the Christian church consists of free, independent, and self-governing congregations” (35).

All this is comprised in the formal principle of Congregationalism:

“The Word of God in Scriptures is designed to furnish, and actually does furnish, the sole objective authority, not only for the doctrines, but also for the constitution, worship, and discipline, of the Christian church.” (Ladd 38).

“It is a principle of congregationalism, that the Scriptures are the only infallible guide in matters of church order and discipline.” (Punchard 19).

Or, as the Declaration of the English Congregationalists of 1833 says:

“They believe that the New Testament contains, either in the form of express statute, or in the example and practice of apostles and apostolic churches, all the articles of faith necessary to be believed, and all the principles of order and discipline requisite for constituting and governing Christian societies; and that human traditions fathers and councils, canons and creeds, possess no authority over the faith and practice of Christians.”

We now understand why S., in a newspaper discussion with Dr. Kildahl, could say: “God and His Son and His Spirit are entirely omitted from Kildahl’s consideration of church polity” (Veiledning 159).

For S. and O. the question of church polity was a fundamental question: church polity had to be derived from Scripture just as truly as doctrine. Not so for Schmidt and Kildahl. And not so for the great majority in the Lutheran church, as is illuminated in the following citation from one of the newest works in theology. After having set forth the Catholic church’s conception of divine right, the author says:

“According to the evangelical conception there is no divine right (church law). The Bible is a source only for faith, not for law. The Symbols are confessional writings, not law-books. For Luther, Christ was no ‘jurist or ruler in outward matters.’ The Lutheran confessional writings do indeed contain a series of thoughts that have become determinative for later legal formations, yet they are nevertheless not sources of law.

\begin{quote}
Somewhat otherwise it is with the Reformed church, insofar as it has seen itself occasioned to adopt the apostolic congregational order. But it too rejects the Roman juridical conception.
\end{quote}

We have now made it clear that principle 1 is a constitutional principle: it is not a dogmatic principle. As a principle of church order, it belongs to ethics, and as applied ethics it belongs to church polity. If one wants the dogmatic side of the congregation, one can find it treated, for example, in Sverdrup’s “Dogmatic paragraphs on the congregation” (S. S. VI.)

Dogmatics sets forth chiefly what God does, His work; ethics, what the human being does, his activity.

Dogmatics presents the congregation as the work of the Spirit, furnished with means of grace and gifts of grace.

Ethics presents the congregation as a gathering of people who know the conditions of salvation, proclaim the Kingdom of God, and labor for its advancement among themselves and others. Therefore it also treats the church’s ministry, organization, relation to the state, tasks, and so forth.

Dogmatics presents the congregation as a gift from God; ethics presents it as a task: God’s gifts become our tasks; He gives the power to carry them out, He uses us as His instruments.

Ethics also treats church polity, but more from an abstract standpoint. A parallel discipline to ethics is practical theology; it applies theological knowledge to practical life, as we live it in experience. Church polity is one of the special branches of practical theology and treats—both theoretically and practically, historically and systematically—the area peculiar to it, which ethics, so to speak, only brushes upon.

Guiding Principles are throughout principles (or fundamental propositions) for church polity. They are principles, not a task. They do not pertain to the individual person, but to the community. They are not an ideal. They do not set forth an unrealizable demand.

Gynild indeed says: “If the same principles did not demand more than what we could carry out, then they would not be Christian.”

An astonishing statement about what Christian principles demand..

Thus perhaps Søren Kierkegaard might have spoken—the spokesman of paradox and of life-denying asceticism, he who called himself a police talent and an angel of accusation, and thought he was appointed to point out what Christianity required and to maintain that there was no one who could fulfill its demands. For K. Christianity was something so exacting that he did not profess himself to be a Christian.

Thus perhaps Bjørnstjerne Bjørnson too might have spoken, when he was working on his drama Over Evne, where—in another spirit than Kierkegaard’s—he sought to show that Christianity is beyond one’s power: its demands cannot be met.

But thus it ought not to be spoken where one wishes to preach evangelical Lutheran Christianity, even if one makes a principle identical with a demand or an ideal.

The highest ideals we possess come from Christianity. The ideals of the ancient world were poor; it did not take much to realize them. Its “god” could not help; for he was so infinitely far away. But the ideals of Christianity lie in the distance. Much is required to realize them. But then the Christian’s God is always near. He gives the power to realize even the highest Christian ideal.

The further away the ideals are, the more exertion is needed in order to attain them (activity). The nearer God is, the more receptivity there is for Him and His gifts (receptivity). He gives us what He requires. And the highest thing He requires is evangelical perfection (Matt. 5:48, not the perfection demanded by the Law). He gives us this perfection (Phil. 3:14, 15: “Let us then, as many as be perfect...”)

The doctrine of Scripture concerning this highest goal or ideal we also find in the Augsburg Confession, art. 27, 12:

\begin{quote}
Christian perfection consists in this: earnestly to fear God, and yet at the same time to set one’s whole confidence in Him, and on account of Christ to trust that we have a reconciled God; to pray to God for help and confidently expect it from Him in all that, according to one’s calling, one has to do; and besides this, in one’s outward conduct, unweariedly to do good works and take heed to one’s calling. In this consists true perfection and true worship of God.
\end{quote}

By Christian perfection is meant neither sinlessness nor the state of being complete. But according to Scripture the predicate perfect belongs to the one who has the moral quality that God wills him to have. The Christian strives to make faith and love a settled part of his spiritual life, so that they determine what he does and what he is.

One of the great thinkers in the realm of Christian thought has said: “The highest ideal of life is Christian perfection. It is a goal that is not of the world to come, but which can be understood and realized here on earth.”

It is precisely the Christian ideals that can be realized.

There is also a heavenly goal. It is not thought to be realized on earth. But it shall be realized there where it is appointed to be realized, when we “see Him as He is” and when the Son “delivers the kingdom to God the Father” (1 Cor. 15:24). There is another order there than in the Christian economy on earth, where we may say: By the grace of God I am what I am (1 Cor. 15:10); I can do all things in Christ who makes me strong. Or: by the grace of God it is possible to keep one’s baptismal covenant.

He who is a Christian must fight. Luther said, “A Christian is not one who has already become one (nicht im Wordensein), but one who is becoming (im Werden) a Christian. He who is a Christian is no Christian.” Or, as Chrysostom said: “To be perfect is: not to think oneself perfect.” Or, as others have expressed it: “The Christian’s perfection consists in this, that he is always conscious of his imperfection.” The more the Christian life increases, the more one comes to know one’s own imperfection; but the more one also learns to know grace, that God does not require more than He gives.

It is possible that this sounds strange to our ears, but it is nothing new—at any rate not new for a man such as, for example, Gustav Jensen, who in his book Fra Menighedsforsamlingen says of Matt. 5:43–48, Be ye perfect:

\begin{quote}
He (Jesus) does not hesitate to set before them the highest goal that it is possible at all to set... A father does not lay before his children what is impossible and unattainable, and when your heavenly Father wills that His children should strive after the highest, then it is because He, as Father, wills to help them attain it. Therefore we also see in the New Testament that the Lord Himself and His apostles often speak of perfection, and indeed of a perfection which already here, according to its basic features, can be attained by His Holy Spirit... Be grasped by Christ, and then be among those who strive after it... And you shall surely advance in that perfection which can be attained here on earth under the grace of the forgiveness of sins, and when the time comes, win the prize.
\end{quote}

This therefore applies to an ideal, a goal (S. would presumably say: a principle) which Jesus Himself has set forth as “the highest goal that it is possible at all to set.” Leading Lutheran men proclaim that it is attainable here on earth. But then Gynild comes and proclaims that Guiding Principles (which are not at all ideals of life and which were, moreover, written by sons of the nineteenth century) are Christian in so high a degree that they demand more than we can carry out: “If the same principles did not demand more than what we could carry out, then they would not be Christian.”

If some solitary schoolmaster had concocted something of that sort, I should have remained silent. But when a pastor, who time and again has been chairman of the Free Church, publishes something like this, then it is right to go against such an assertion.

Guiding Principles are not demands, are not ideals. They are principles.

What is a principle?

The word principle is used in very different senses. In philosophy—which, to use a Gynildian word, the “learned fools” must now and then condescend to concern themselves a little with—principle means: highest concept, supreme proposition, law, axiom, explanatory ground, point of departure. A favorite definition of philosophy is the science of principles.

In the science of religion, principle means the conception of the fundamental impulse or fundamental power that lies behind the individual psychological expressions of life and facts.

In historical psychology, “principle” is applied to a totality determined by certain interests, for example in the analysis of an individual, the character of a people, a section of culture. Here principle can be paraphrased by words such as national spirit, spirit of the age, essence, fundamental view, fundamental idea, idea, and so forth.

The psychological conception of religion sees in “principle” not ready-made propositions, but a fundamental disposition of the soul: the principle is the point of unity behind consciousness.

